\documentclass[a4paper,11pt]{article}
\usepackage[utf8]{inputenc}
\usepackage[T1]{fontenc}
\usepackage[french]{babel}
\usepackage{lmodern}
\usepackage{amsmath,amssymb,amsthm}
\usepackage{mathrsfs}
\usepackage{enumitem}
\usepackage{multicol}

\setlist[enumerate]{label=\textbf{\textcolor{Couleur8}{\arabic*.}}, left=1em}

% Si vous voulez aussi modifier les listes enumerate imbriquées :
\setlist[enumerate,2]{label=\textcolor{Couleur8}{\alph*)}}
\setlist[enumerate,3]{label=\textcolor{Couleur8}{\roman*)}}
\setlist[enumerate,4]{label=\textcolor{Couleur8}{\Alph*)}}
\usepackage{graphicx}
\usepackage{xcolor}
\usepackage{tcolorbox}
\usepackage{geometry}
\usepackage{fancyhdr}
\usepackage{titlesec}
\usepackage{cancel}
\usepackage{ifthen}
\usepackage{tikz}
\usetikzlibrary{arrows.meta}
\usepackage{tkz-tab}
\usepackage{eurosym}

% OPTION PROF/ÉLÈVE
% Décommentez UNE des deux lignes suivantes :
\newboolean{prof}\setboolean{prof}{true}   % Version professeur (avec corrections des devoirs)
\newboolean{prof}\setboolean{prof}{true}  % Version élève (sans corrections des devoirs)

\definecolor{Couleur1}{HTML}{4A5D7A} %Th
\definecolor{Couleur2}{HTML}{A5B5D6} %Prop
\definecolor{Couleur3}{HTML}{A8C68A} %Méthode
\definecolor{Couleur6}{HTML}{8A9CC1} %Def
\definecolor{Couleur7}{HTML}{E6B459} %Rq Voc
\definecolor{Couleur8}{HTML}{D36740} %Titres
\definecolor{correction}{HTML}{D36740} % Couleur pour les corrections
\definecolor{coopmathscorr}{HTML}{F15929} % Couleur corrections coopmaths

% Configuration des marges
\geometry{
	top=2cm,
	bottom=2cm,
	inner=1.5cm,
	outer=1.5cm,
}

% Police
\renewcommand{\familydefault}{\sfdefault}

% Compteurs
\newcounter{seancenum}
\newcounter{seancenumD}
\setcounter{seancenum}{1}
\setcounter{seancenumD}{1}

% Configuration des en-têtes et pieds de page
\pagestyle{fancy}
\fancyhf{}
\ifthenelse{\boolean{prof}}{
    \fancyhead[L]{\textcolor{Couleur8}{\textbf{Corrections Automatismes - Classe de seconde - Version Professeur}}}
}{
    \fancyhead[L]{\textcolor{Couleur8}{\textbf{Corrections Devoirs - Séance 9 - Classe de seconde}}}
}
\fancyhead[R]{\textcolor{Couleur8}{\textbf{2025-2026}}}
\fancyfoot[L]{\textcolor{Couleur8}{Mme Bertail et Mme Pinto}}
\fancyfoot[R]{\textcolor{Couleur8}{\thepage}}
\renewcommand{\headrulewidth}{1pt}
\renewcommand{\headrule}{\hbox to\headwidth{\color{Couleur8}\leaders\hrule height \headrulewidth\hfill}}

% Environnements personnalisés
\newtcolorbox{seance}[1]{
    colback=Couleur6!10,
    colframe=Couleur1!65,
    fonttitle=\bfseries\large,
    title=Correction Séance \theseancenum #1,
    left=5mm,
    right=5mm,
    top=3mm,
    bottom=3mm,
    arc=0mm,
    after=\stepcounter{seancenum}
}

\newtcolorbox{seanceD}[1]{
    colback=Couleur6!5,
    colframe=Couleur6!45,
    fonttitle=\bfseries\large,
    title=Correction Devoirs - Séance \theseancenumD #1,
    left=5mm,
    right=5mm,
    top=3mm,
    bottom=3mm,
    arc=0mm,
    after=\stepcounter{seancenumD}
}

\begin{document}

\setcounter{seancenumD}{9}
\newpage
\begin{seanceD}{} %9
\begin{enumerate}
\item Diminuer de $80\,\%$ revient à multiplier par $0{,}2$ et augmenter de $10\,\%$ revient à multiplier par $1{,}1$.\\
  Globalement cela revient donc à multiplier par $0{,}2\times 1{,}1=0{,}22$.\\
  Multiplier par $0{,}22$ revient à multiplier par  $1-0{,}78$. \\
          Le taux d'évolution global est donc : ${\color[HTML]{f15929}\boldsymbol{-78\,}} \%$.
          
\item  Le coefficient multiplicateur associé à une baisse de $50\,\%$ est $1-0{,}5=0{,}5$.\\
    Comme $0{,}5\times 2=1$, il faut donc multiplier par $2$ (coefficient multiplicateur réciproque) pour revenir au prix initial.\\
      Un coefficient multiplicateur de $2$ correspond à un taux d'évolution de $100\,\%$.\\
      On en déduit qu'il faut une augmentation de ${\color[HTML]{f15929}\boldsymbol{100\,\%}}$.\\La bonne réponse est la réponse {\bfseries \color[HTML]{f15929}D}.


\item Il s'agit d'un problème additif. Il va être nécessaire de réduire les fractions au même dénominateur pour les additionner, les soustraire ou les comparer.\\Réduisons les fractions de l'énoncé au même dénominateur :  $\dfrac{11}{40}$  et $\dfrac{1}{4}$ $= \dfrac{10}{40}$.\\Calculons d'abord la distance en course à pied : \\$1-\dfrac{11}{40}-\dfrac{1}{4} = \dfrac{40}{40}-\dfrac{11}{40}-\dfrac{10}{40} = \dfrac{40-11-10}{40} = \dfrac{19}{40}$\\Gaspard fait donc $\dfrac{11}{40}$ en VTT, $\dfrac{1}{4}$ en ski de fond et $\dfrac{19}{40}$ en course à pied.\\ Avec les mêmes dénominateurs pour pouvoir comparer, Gaspard fait donc $\dfrac{11}{40}$ en VTT, $\dfrac{10}{40}$ en ski de fond et $\dfrac{19}{40}$ en course à pied.\\Nous pouvons alors ranger ces fractions dans l'ordre croissant : $\dfrac{10}{40}$ < $\dfrac{11}{40}$ < $\dfrac{19}{40}$.\\Enfin, nous pouvons ranger les fractions de l'énoncé et la fraction calculée dans l'ordre croissant : ${\color[HTML]{f15929}\boldsymbol{\dfrac{1}{4}}}$ < ${\color[HTML]{f15929}\boldsymbol{\dfrac{11}{40}}}$ < ${\color[HTML]{f15929}\boldsymbol{\dfrac{19}{40}}}$.\\ 
                      C'est donc en {\bfseries \color[HTML]{f15929}course à pied} que Gaspard fait la plus grande distance.

\item Pour $x\in \mathbb{R}\smallsetminus\left\{-1\right\}$, \\
            $\begin{aligned}
            4x+5+\dfrac{2}{3x+3}
           & = \dfrac{(4x+5)(3x+3)}{3x+3}+\dfrac{2}{3x+3}\\
            &=\dfrac{(12x^2+27x+15)+2}{3x+3}\\
           & =\dfrac{12x^2+27x+17}{3x+3}
           \end{aligned}$
\item $2-(-6x-6)=x-5$\\On développe le membre de gauche.\\$2+6x+6=x-5$\\$6x+8=x-5$\\On soustrait $x$ aux deux membres.\\$6x+8{\color[HTML]{216D9A}\boldsymbol{-x}}=1x+-5{\color[HTML]{216D9A}\boldsymbol{-x}}$\\$5x+8=-5$\\On soustrait $8$ aux deux membres.\\$5x+8{\color[HTML]{216D9A}\boldsymbol{-8}}=-5{\color[HTML]{216D9A}\boldsymbol{-8}}$\\$5x=-13$\\On divise les deux membres par $5$.\\$5x{\color[HTML]{216D9A}\boldsymbol{\div5}}=-13{\color[HTML]{216D9A}\boldsymbol{\div5}}$\\$x=-\dfrac{13}{5}$\\La solution de l'équation $2-(-6x-6)=x-5$ est ${\color[HTML]{f15929}\boldsymbol{-\dfrac{13}{5}}}$.
\end{enumerate}
\end{seanceD}

\end{document}